\section{Introduction}
\label{sec:intro}
Leaderboards report each model on dozens to hundreds of benchmarks, and running a new model on all of them is expensive. Yet the scores of earlier models are already known, and benchmark scores of different models are strongly correlated: the score matrices studied here have a leading component that carries $36$--$80\%$ of the variance. A natural practice is therefore to run a few benchmarks and infer the rest. Work on cheaper evaluation has mostly reduced the number of \emph{items} inside one benchmark \citep{polo2024tiny,kipnis2024metabench}, or fitted low-dimensional capability models to scores of many models \citep{ruan2024observational,polo2024sloth}. We ask the complementary question at the level of whole benchmarks, in the setting an evaluator faces for a model that has nothing in the table yet:
\begin{quote}\emph{Given a historical score matrix with missing cells, and a new model that can run only $k$ benchmarks: which $k$ should it run, how should the remaining scores be inferred, and how large must $k$ be?}\end{quote}
The three parts interact, and ``how large'' has no meaning without a yardstick for ``enough''. We therefore fix the evaluation before looking at final results: the unit of evaluation is a held-out model; the selection order and every fitted quantity use training models only; the targets are a fixed set of benchmarks that no selector may pick, so every $k$ is judged on the same cells; uncertainty is a bootstrap over held-out models; and $k$ is called enough when it closes $80\%$ or $90\%$ of the gap between the column mean and the full-information predictor of the same inference rule (lower confidence bound), with absolute tolerances as a secondary criterion.

\paragraph{Contributions.}
\begin{itemize}\itemsep0pt
\item A cold-start protocol and a pre-fixed criterion for the size of $k$ (Sections~\ref{sec:method}, \ref{sec:protocol}), applied to nine real matrices (four MTEB blocks of embedding models, one of them $30\%$ observed; generative LLMs on HELM-MMLU, two HELM Lite matrices, LiveBench and the Open LLM Leaderboard) and to synthetic matrices with known truth.
\item An empirical comparison of $10$ selectors and $7$ inference rules. Informed selection that accounts for redundancy (greedy conditional variance, greedy log-determinant) beats random selection by $0.08$--$0.15$ nMSE at $k\le4$ on MTEB-56; common column-subset heuristics (leverage scores, pivoted QR, highest variance) do not, and a one-line centrality heuristic matches them only at $k\le2$. Direct contrasts show that no informed selector is uniformly best and the advantage is $\le0.03$ nMSE by $k\approx16$ (Section~\ref{sec:exp}).
\item Answers to ``how large $k$'': a handful closes most of the gap ($k=2$--$4$ for $80\%$ and $k=4$--$16$ for $90\%$ on MTEB-56, informed selection) to the full-information predictor, but that predictor is itself far from perfect ($\nmse=0.15$--$0.36$), and approaching it takes tens of benchmarks; the bootstrap interval of $k^\star$ and its value under leave-organisation-out and leave-family-out splits are reported, and the ranking of informed selectors is not stable across them.
\item Stress tests (missing cells, few training models, rank, leave-family-out and temporal shift, leakage, noise, cost), with failures reported (Sections~\ref{sec:exp}, \ref{sec:disc}).
\item Bounds for random subsets and a $(1-1/e)$ guarantee for greedy information selection, checked against synthetic and real matrices (Appendix~\ref{app:theory}, Fig.~\ref{fig:theory}). They are a sanity check on the form of the $k$-curve, not a contribution about its level on real data.
\end{itemize}
All numbers come from the code shipped with the paper (Appendix~\ref{app:repro}).
